\documentclass[11pt]{article}

\usepackage[margin=1in]{geometry}
\usepackage{amsmath}
\usepackage{amssymb}
\usepackage{booktabs}
\usepackage{enumitem}
\usepackage[hidelinks]{hyperref}
\usepackage{mathtools}

\newcommand{\R}{\mathbb{R}}
\newcommand{\norm}[1]{\left\lVert #1 \right\rVert}
\newcommand{\abs}[1]{\left\lvert #1 \right\rvert}
\newcommand{\set}[1]{\left\{ #1 \right\}}
\newcommand{\trans}{^{\mathsf T}}
\newcommand{\level}{\ell}
\DeclareMathOperator{\diag}{diag}
\DeclareMathOperator{\sgn}{sgn}
\DeclareMathOperator{\nnz}{nnz}

\title{A Minimal PMIS--Extended+i--L1-Jacobi\\
       Algebraic Multigrid Method for AMReX}
\author{AMReX Design Note}
\date{\today}

\begin{document}

\maketitle

\begin{abstract}
This note specifies the equations for a minimal, classical C/F algebraic
multigrid solver built on AMReX distributed sparse-matrix infrastructure.  The
method uses classical strength of connection with the maximum-row-sum rule,
parallel modified independent set (PMIS) coarsening, matrix-based Extended+i
interpolation capped at four entries per row, transpose restriction, Galerkin
coarse operators, L1-Jacobi smoothing, a multiplicative V(1,1)-cycle, and a
small dense LU coarse solve or a terminal L1-Jacobi sweep when setup stops
early.  It is intended to accompany
\texttt{AMG\_BoomerAMGLite\_Plan.md}.  The comparison profile corresponds to
BoomerAMG coarsening type 8, interpolation type 6, and relaxation type 18; it
does not use aggregation.
\end{abstract}

\tableofcontents

\section{Linear System and Hierarchy}

Let the fine-level problem be
\begin{equation}
  A_0 x_0 = b_0,
  \qquad
  A_0 \in \R^{n_0\times n_0},
  \qquad
  x_0,b_0\in\R^{n_0}.
  \label{eq:fine-system}
\end{equation}
Level indices increase as the problem is coarsened:
\begin{equation}
  n_0 > n_1 > \cdots > n_L,
  \qquad
  A_\level\in\R^{n_\level\times n_\level}.
  \label{eq:level-sizes}
\end{equation}
The first implementation assumes that every $A_\level$ is square, nonsingular,
and has a nonzero diagonal.  Its principal target is a scalar elliptic matrix
with
\begin{equation}
  a_{ii}>0,
  \qquad
  a_{ij}\leq 0 \quad (i\neq j),
  \label{eq:m-matrix-sign}
\end{equation}
although the sign-separated interpolation equations below admit mixed
off-diagonal signs.

On MPI rank $p$, the global row interval is described by a prefix partition
\begin{equation}
  \mathcal I_p =
  \set{r_p,r_p+1,\ldots,r_{p+1}-1},
  \qquad
  0=r_0\leq r_1\leq\cdots\leq r_{N_p}=n_\level.
  \label{eq:row-partition}
\end{equation}
The sparse matrix is split into locally owned and off-rank columns,
\begin{equation}
  A_\level =
  \begin{bmatrix}
    A_\level^{\mathrm{loc}} &
    A_\level^{\mathrm{off}}
  \end{bmatrix},
  \label{eq:parcsr-split}
\end{equation}
with a cached halo exchange supplying vector values addressed by
$A_\level^{\mathrm{off}}$.

\section{Step 1: Classical Strength of Connection}

For each off-diagonal entry in row $i$, define its sign-adjusted coupling
\begin{equation}
  q_{ij} = -\sgn(a_{ii})a_{ij},
  \qquad j\neq i,
  \label{eq:sign-adjusted-coupling}
\end{equation}
and the largest opposite-sign coupling
\begin{equation}
  m_i = \max_{k\neq i}\max(q_{ik},0).
  \label{eq:row-strength-scale}
\end{equation}
Let the diagonal-inclusive row sum be
\begin{equation}
  r_i=\sum_j a_{ij}.
  \label{eq:matrix-row-sum}
\end{equation}
Given a strength threshold $\theta\in(0,1]$ and maximum-row-sum parameter
$\rho\in(0,1]$, row $i$ strongly depends on column $j$ when
\begin{equation}
  S_{ij} =
  \begin{cases}
    1, & q_{ij}\geq\theta m_i,\quad m_i>0,\quad j\neq i,
         \quad \abs{r_i}\leq\rho\abs{a_{ii}},\\
    0, & \text{otherwise}.
  \end{cases}
  \label{eq:strength-definition}
\end{equation}
The default profile uses
\begin{equation}
  \theta=0.25,
  \qquad
  \rho=0.9.
  \label{eq:strength-threshold}
\end{equation}
Thus every off-diagonal entry is weak when
$\abs{r_i}>0.9\abs{a_{ii}}$.  The implementation treats $\rho=1$ as disabling
this all-weak-row rule.
For the positive-diagonal M-matrix case in
\eqref{eq:m-matrix-sign}, this reduces to
\begin{equation}
  -a_{ij}
  \geq
  \theta\max_{k\neq i}(-a_{ik}).
  \label{eq:m-matrix-strength}
\end{equation}

Only indices satisfying $S_{ij}=1$ are stored.  Thus the strength graph has no
numerical zero entries and no diagonal:
\begin{equation}
  \nnz(S)
  =
  \abs{\set{(i,j):S_{ij}=1}},
  \qquad
  S_{ii}=0.
  \label{eq:compact-strength}
\end{equation}

\section{Step 2: PMIS Priorities}

The influence measure for point $i$ is the number of equations that strongly
depend on it:
\begin{equation}
  \mu_i
  =
  \sum_k S_{ki}
  =
  (S\trans\mathbf 1)_i.
  \label{eq:influence-measure}
\end{equation}
This is obtained in the prototype design by constructing $S\trans$ and taking
its row sum.

Let $g_i$ be the global row ID and let
\begin{equation}
  h(g_i,s)\in\set{0,1,\ldots,2^{64}-1}
  \label{eq:hash-range}
\end{equation}
be a deterministic hash with seed $s$.  Convert it to a fractional tie-breaker
\begin{equation}
  u_i
  =
  \frac{h(g_i,s)+\tfrac12}{2^{64}},
  \qquad
  0<u_i<1,
  \label{eq:priority-fraction}
\end{equation}
and define the PMIS priority
\begin{equation}
  \pi_i=\mu_i+u_i.
  \label{eq:pmis-priority}
\end{equation}
The integer part favors highly influential variables, while the fractional
part establishes a deterministic total ordering.  A final lexicographic
comparison with $g_i$ resolves the exceptional case of a hash collision:
\begin{equation}
  (\pi_i,g_i)>_{\mathrm{lex}}(\pi_j,g_j).
  \label{eq:lexicographic-priority}
\end{equation}

\section{Step 3: PMIS C/F Splitting}

Let $\mathcal U^{(t)}$ be the undecided set at PMIS round $t$.  Define a
conflict neighborhood from the directed strength graph by
\begin{equation}
  \mathcal N_i^{\mathrm{conf}}
  =
  \set{j:S_{ij}=1\ \text{or}\ S_{ji}=1}.
  \label{eq:conflict-neighborhood}
\end{equation}
The independent coarse candidates in round $t$ are
\begin{equation}
  \mathcal C^{(t)}
  =
  \set{
    i\in\mathcal U^{(t)}:
    (\pi_i,g_i)>_{\mathrm{lex}}(\pi_j,g_j)
    \ \forall j\in
    \mathcal N_i^{\mathrm{conf}}\cap\mathcal U^{(t)}
  }.
  \label{eq:round-coarse-set}
\end{equation}
An undecided point becomes fine if it strongly depends on a newly selected
coarse point:
\begin{equation}
  \mathcal F^{(t)}
  =
  \set{
    i\in\mathcal U^{(t)}\setminus\mathcal C^{(t)}:
    \exists c\in\mathcal C^{(t)}
    \text{ such that }S_{ic}=1
  }.
  \label{eq:round-fine-set}
\end{equation}
The active set update is
\begin{equation}
  \mathcal U^{(t+1)}
  =
  \mathcal U^{(t)}
  \setminus
  \left(\mathcal C^{(t)}\cup\mathcal F^{(t)}\right).
  \label{eq:active-set-update}
\end{equation}
The process stops when the global undecided count is zero:
\begin{equation}
  \sum_{p=0}^{N_p-1}
  \abs{\mathcal U_p^{(t)}}=0.
  \label{eq:pmis-termination}
\end{equation}

The final splitting is
\begin{equation}
  \mathcal C=\bigcup_t\mathcal C^{(t)},
  \qquad
  \mathcal F=\bigcup_t\mathcal F^{(t)},
  \qquad
  \mathcal C\cap\mathcal F=\varnothing,
  \qquad
  \mathcal C\cup\mathcal F=\set{0,\ldots,n_\level-1}.
  \label{eq:final-splitting}
\end{equation}
Isolated points are promoted to $\mathcal C$.  No aggregate mapping is
constructed.

The per-round independent-set invariant and the final interpolation-coverage
invariant are
\begin{align}
  &\text{round independence:}
  & i,j\in\mathcal C^{(t)},\ i\neq j
  &\Longrightarrow
  j\notin\mathcal N_i^{\mathrm{conf}},
  \label{eq:round-coarse-independence}\\
  &\text{interpolation coverage:}
  & i\in\mathcal F
  &\Longrightarrow
  \exists c\in\mathcal C\text{ with }S_{ic}=1.
  \label{eq:fine-coverage}
\end{align}
Because fine points are removed only when they strongly depend on a new coarse
point, a later PMIS round can select a point joined to an earlier coarse point
by an edge in the opposite direction.  The final coarse set therefore need
not be independent in the symmetrized strength graph.

\section{Step 4: Distributed Coarse Numbering}

Let rank $p$ own $n_{\level,p}$ fine rows.  Its local number of coarse points
is
\begin{equation}
  n_{\level+1,p}
  =
  \sum_{i\in\mathcal I_p}
  \chi_{\mathcal C}(i),
  \label{eq:local-coarse-count}
\end{equation}
where $\chi_{\mathcal C}$ is the indicator of the coarse set.  The rank's
global coarse offset is
\begin{equation}
  o_{\level+1,p}
  =
  \sum_{q=0}^{p-1}n_{\level+1,q},
  \label{eq:coarse-rank-offset}
\end{equation}
and the next global problem size is
\begin{equation}
  n_{\level+1}
  =
  \sum_{p=0}^{N_p-1}n_{\level+1,p}.
  \label{eq:global-coarse-count}
\end{equation}

For a locally owned coarse point $i$, its global coarse ID is
\begin{equation}
  I_c(i)
  =
  o_{\level+1,p}
  +
  \sum_{\substack{k\in\mathcal I_p\\k<i}}
  \chi_{\mathcal C}(k).
  \label{eq:coarse-id}
\end{equation}
The second term is a device exclusive scan.  Values $I_c(i)$ needed by
off-rank interpolation rows are exchanged as integer fields.

\section{Step 5: Matrix-Based Extended+i Interpolation}

The interpolation operator has dimensions
\begin{equation}
  P_\level\in\R^{n_\level\times n_{\level+1}}.
  \label{eq:p-dimensions}
\end{equation}
For a coarse point, interpolation is injection:
\begin{equation}
  (P_\level)_{i,I_c(i)}=1,
  \qquad
  (P_\level)_{i,k}=0\quad(k\neq I_c(i)),
  \qquad i\in\mathcal C.
  \label{eq:coarse-interpolation-row}
\end{equation}

For a fine point, define its strong coarse and strong fine neighborhoods
\begin{equation}
  \mathcal C_i^s
  =
  \set{j\in\mathcal C:S_{ij}=1},
  \qquad
  \mathcal F_i^s
  =
  \set{j\in\mathcal F:S_{ij}=1}.
  \label{eq:strong-coarse-neighborhood}
\end{equation}
The diagonal and weak connections contribute
\begin{equation}
  w_i
  =
  a_{ii}
  +
  \sum_{\substack{j\neq i\\S_{ij}=0}}a_{ij},
  \label{eq:weak-diagonal-sum}
\end{equation}
while the strong-coarse row sum is
\begin{equation}
  \gamma_i
  =
  \sum_{c\in\mathcal C_i^s}a_{ic}.
  \label{eq:strong-coarse-sum}
\end{equation}
For a strong fine connection $j\in\mathcal F_i^s$, define
\begin{equation}
  \delta_{ij}=\gamma_j+a_{ji},
  \qquad
  q_{ij}
  =
  \begin{cases}
    a_{ij}/\delta_{ij}, & \delta_{ij}\neq0,\\
    0,                  & \delta_{ij}=0.
  \end{cases}
  \label{eq:extended-i-fine-scale}
\end{equation}
The Extended+i row denominator is
\begin{equation}
  \eta_i
  =
  w_i
  +
  \sum_{j\in\mathcal F_i^s}
  \begin{cases}
    a_{ji}q_{ij}, & \delta_{ij}\neq0,\\
    a_{ij},       & \delta_{ij}=0.
  \end{cases}
  \label{eq:extended-i-row-denominator}
\end{equation}
Define a scaled fine-fine matrix $M_{FF}$ by
\begin{equation}
  (M_{FF})_{ij}
  =
  \begin{cases}
    -1/\eta_i,      & j=i,\\
    -q_{ij}/\eta_i, & j\in\mathcal F_i^s,\\
    0,              & \text{otherwise},
  \end{cases}
  \qquad i\in\mathcal F.
  \label{eq:scaled-ff}
\end{equation}
If $A_{FC}^s$ contains the strong fine-to-coarse entries, the fine block of
the untruncated interpolation operator is the sparse product
\begin{equation}
  P_F=M_{FF}A_{FC}^s.
  \label{eq:extended-i-product}
\end{equation}
Coarse rows are represented by identity rows in both factors, so the product
also gives the injection rows in \eqref{eq:coarse-interpolation-row}.

After forming the product, let $K_i$ contain at most the four entries of row
$i$ with largest magnitude, breaking equal-magnitude ties by global coarse
column ID, and let $j_i^*$ be the first retained column in sorted CSR order.
With
\begin{equation}
  s_i=\sum_j(P_\level)_{ij},
  \qquad
  \widehat s_i=\sum_{j\in K_i}(P_\level)_{ij},
  \label{eq:truncation-row-sums}
\end{equation}
the stored interpolation entries are
\begin{equation}
  (\widehat P_\level)_{ij}
  =
  \begin{cases}
    \dfrac{s_i}{\widehat s_i}(P_\level)_{ij},
      & j\in K_i,\quad \widehat s_i\neq0,\\[2mm]
    (P_\level)_{ij}+s_i,
      & j=j_i^*,\quad \widehat s_i=0,\\
    (P_\level)_{ij},
      & j\in K_i\setminus\set{j_i^*},\quad \widehat s_i=0,\\
    0, & j\notin K_i.
  \end{cases}
  \label{eq:interpolation-truncation}
\end{equation}
Thus truncation preserves the original row sum.  If the retained entries
cancel exactly but $s_i\neq0$, the implementation adds $s_i$ to the first
retained entry.  A configured cap of zero disables truncation.

\section{Step 6: Restriction and Galerkin Coarsening}

Restriction is the transpose of interpolation:
\begin{equation}
  R_\level=P_\level\trans
  \in\R^{n_{\level+1}\times n_\level}.
  \label{eq:transpose-restriction}
\end{equation}
The coarse operator is Galerkin:
\begin{equation}
  A_{\level+1}
  =
  R_\level A_\level P_\level
  =
  P_\level\trans A_\level P_\level.
  \label{eq:galerkin-operator}
\end{equation}
It is formed with two distributed sparse products:
\begin{equation}
  B_\level=A_\level P_\level,
  \qquad
  A_{\level+1}=R_\level B_\level.
  \label{eq:two-stage-rap}
\end{equation}
The action identity used by tests is
\begin{equation}
  A_{\level+1}v
  =
  R_\level\left(A_\level(P_\level v)\right)
  \qquad
  \forall v\in\R^{n_{\level+1}}.
  \label{eq:galerkin-action-test}
\end{equation}
If $A_\level$ is symmetric positive definite and $P_\level$ has full column
rank, then
\begin{equation}
  v\trans A_{\level+1}v
  =
  (P_\level v)\trans A_\level(P_\level v)>0
  \qquad(v\neq0),
  \label{eq:spd-preservation}
\end{equation}
so Galerkin coarsening preserves positive definiteness.

\section{Step 7: L1-Jacobi Smoothing}

For every level, cache the row L1 denominator
\begin{equation}
  d_{\level,i}
  =
  \sum_j\abs{(A_\level)_{ij}},
  \qquad
  D_\level^{(1)}
  =
  \diag(d_{\level,0},\ldots,d_{\level,n_\level-1}).
  \label{eq:l1-diagonal}
\end{equation}
One L1-Jacobi sweep with weight $\omega=1$ is
\begin{equation}
  \mathcal S_\level(x,b)
  =
  x+
  \left(D_\level^{(1)}\right)^{-1}
  (b-A_\level x).
  \label{eq:l1-jacobi}
\end{equation}
Componentwise,
\begin{equation}
  x_i^{\mathrm{new}}
  =
  x_i+
  \frac{b_i-(A_\level x)_i}
       {\sum_j\abs{(A_\level)_{ij}}}.
  \label{eq:l1-jacobi-component}
\end{equation}
For a zero initial correction, the first sweep simplifies to
\begin{equation}
  \mathcal S_\level(0,b)
  =
  \left(D_\level^{(1)}\right)^{-1}b,
  \label{eq:zero-guess-smoother}
\end{equation}
and no sparse matrix-vector product is needed.

\section{Step 8: Multiplicative V(1,1)-Cycle}

Define
\begin{equation}
  \mathcal V_\level(b_\level,x_\level)
  \label{eq:vcycle-operator}
\end{equation}
as one V-cycle on level $\level$.  For $\level<L$, the pre-smoothing step is
\begin{equation}
  \bar x_\level
  =
  \mathcal S_\level(x_\level,b_\level).
  \label{eq:pre-smoothing}
\end{equation}
The fine residual is
\begin{equation}
  r_\level
  =
  b_\level-A_\level\bar x_\level.
  \label{eq:level-residual}
\end{equation}
The coarse error equation receives the restricted right-hand side
\begin{equation}
  b_{\level+1}
  =
  R_\level r_\level,
  \qquad
  e_{\level+1}^{(0)}=0.
  \label{eq:residual-restriction}
\end{equation}
It is solved recursively:
\begin{equation}
  e_{\level+1}
  =
  \mathcal V_{\level+1}
  (b_{\level+1},e_{\level+1}^{(0)}).
  \label{eq:recursive-coarse-solve}
\end{equation}
The prolongated correction is
\begin{equation}
  \tilde x_\level
  =
  \bar x_\level+P_\level e_{\level+1}.
  \label{eq:coarse-correction}
\end{equation}
The cycle finishes with one post-sweep:
\begin{equation}
  \mathcal V_\level(b_\level,x_\level)
  =
  \mathcal S_\level(\tilde x_\level,b_\level).
  \label{eq:post-smoothing}
\end{equation}
Equations \eqref{eq:pre-smoothing} and \eqref{eq:post-smoothing} give one sweep
in each direction, hence a V(1,1)-cycle.

\section{Step 9: Terminal-Level Solve}

Coarsening normally stops when
\begin{equation}
  n_L\leq 9.
  \label{eq:coarse-size}
\end{equation}
Embed the active coarse matrix in a fixed $9\times9$ matrix:
\begin{equation}
  \widetilde A_L
  =
  \begin{bmatrix}
    A_L & 0\\
    0   & I_{9-n_L}
  \end{bmatrix},
  \qquad
  \widetilde b_L
  =
  \begin{bmatrix}
    b_L\\
    0
  \end{bmatrix}.
  \label{eq:padded-coarse-system}
\end{equation}
A pivoted LU factorization is computed once during setup:
\begin{equation}
  \Pi\widetilde A_L=LU.
  \label{eq:coarse-lu}
\end{equation}
Each V-cycle solves
\begin{align}
  Ly &= \Pi\widetilde b_L,
  \label{eq:coarse-forward}\\
  U\widetilde e_L &= y,
  \label{eq:coarse-backward}
\end{align}
and retains the first $n_L$ entries:
\begin{equation}
  e_L
  =
  \left(\widetilde e_L\right)_{0:n_L-1}.
  \label{eq:coarse-active-solution}
\end{equation}
The identity padding makes this equivalent to solving $A_Le_L=b_L$.

If PMIS does not reduce a level larger than nine rows, or the 25-level limit
is reached first, setup terminates without dense factorization.  Starting from
the V-cycle's zero coarse correction, the terminal action is instead
\begin{equation}
  e_L=\left(D_L^{(1)}\right)^{-1}b_L,
  \label{eq:terminal-l1-jacobi}
\end{equation}
which is one L1-Jacobi sweep and matches BoomerAMG's early-termination
behavior.

\section{Step 10: Preconditioner and Stationary Solver}

The preconditioner applies exactly one zero-initialized V-cycle:
\begin{equation}
  z=M^{-1}r
  \coloneqq
  \mathcal V_0(r,0).
  \label{eq:preconditioner-apply}
\end{equation}
There is no convergence test inside $M^{-1}$.  Consequently, the same fixed
operation can be supplied as a right preconditioner to GMRES:
\begin{equation}
  AM^{-1}y=b,
  \qquad
  x=M^{-1}y.
  \label{eq:right-preconditioned-system}
\end{equation}

A convenience standalone AMG solve uses stationary defect correction:
\begin{align}
  r^{(k)}
  &=
  b-Ax^{(k)},
  \label{eq:stationary-residual}\\
  x^{(k+1)}
  &=
  x^{(k)}+\mathcal V_0(r^{(k)},0).
  \label{eq:stationary-update}
\end{align}
The true residual is recomputed after every cycle; an internal recurrence is
not used for the stopping decision.

\section{Step 11: Tolerance and HYPRE Agreement}

Given an initial guess $x^{(0)}$, define
\begin{equation}
  r^{(0)}=b-Ax^{(0)}
  \label{eq:initial-residual}
\end{equation}
and the requested residual threshold
\begin{equation}
  \tau
  =
  \max\left(
    \mathtt{atol},
    \mathtt{rtol}\norm{r^{(0)}}_2
  \right).
  \label{eq:residual-target}
\end{equation}
The initial benchmark uses $x^{(0)}=0$, so
\begin{equation}
  \norm{r^{(0)}}_2=\norm b_2.
  \label{eq:zero-guess-residual}
\end{equation}

Let $x_A$ be the AMReX result and $x_H$ the result from HYPRE BoomerAMG with
the matching PMIS/Extended+i/L1-Jacobi configuration.  Both independently
recomputed true residuals must satisfy
\begin{align}
  \norm{b-Ax_A}_2
  &\leq\tau,
  \label{eq:amrex-residual-pass}\\
  \norm{b-Ax_H}_2
  &\leq\tau.
  \label{eq:hypre-residual-pass}
\end{align}
On the conditioned manufactured benchmark, solution agreement is
\begin{equation}
  \norm{x_A-x_H}_2
  \leq
  \mathtt{atol}
  +
  \mathtt{rtol}\norm{x_H}_2.
  \label{eq:solution-agreement}
\end{equation}
If the manufactured exact solution $x_\star$ is known, also report
\begin{equation}
  E_A=\norm{x_A-x_\star}_2,
  \qquad
  E_H=\norm{x_H-x_\star}_2.
  \label{eq:exact-solution-errors}
\end{equation}

The tests do not require equality of PMIS markers, interpolation entries,
level counts, or iteration counts.  Different deterministic priorities and
sparse accumulation orders can produce different valid hierarchies while
still satisfying \eqref{eq:amrex-residual-pass}--\eqref{eq:solution-agreement}.

\section{Step 12: Complexity Diagnostics}

Correctness is the first acceptance criterion, but setup records the grid
complexity
\begin{equation}
  C_g
  =
  \frac{\displaystyle\sum_{\level=0}^{L}n_\level}{n_0},
  \label{eq:grid-complexity}
\end{equation}
operator complexity
\begin{equation}
  C_A
  =
  \frac{\displaystyle\sum_{\level=0}^{L}\nnz(A_\level)}
       {\nnz(A_0)},
  \label{eq:operator-complexity}
\end{equation}
and interpolation complexity
\begin{equation}
  C_P
  =
  \frac{\displaystyle\sum_{\level=0}^{L-1}\nnz(P_\level)}
       {n_0}.
  \label{eq:interpolation-complexity}
\end{equation}
These quantities expose poor coarsening or interpolation growth before timing
comparisons are interpreted.

\appendix

\section{Matching BoomerAMG Parameters}

\begin{center}
\begin{tabular}{@{}lll@{}}
\toprule
Feature & HYPRE value & Mathematical section\\
\midrule
Strength threshold & $0.25$ & Section 2\\
Maximum row sum & $0.9$ & Section 2\\
PMIS coarsening & type 8 & Sections 3--4\\
Extended+i interpolation & type 6 & Section 6\\
Interpolation cap & \texttt{PMaxElmts = 4} & Section 6\\
No threshold-based dropping & \texttt{TruncFactor = 0} & Section 6\\
L1-Jacobi down/up & type 18 & Section 8\\
Relaxation order & 0 & Section 8\\
Sweeps & one down, one up & Section 9\\
Transpose restriction & default & Section 7\\
V-cycle & type 1 & Section 9\\
Maximum coarse size & 9 & Section 10\\
Coarse Gaussian elimination & type 9 & Section 10\\
Aggregation levels & 0 & all sections\\
Preconditioner iterations & 1 & Section 11\\
Preconditioner tolerance & 0 & Section 11\\
\bottomrule
\end{tabular}
\end{center}

The supported GPU options and their public integer identifiers are documented
at
\url{https://hypre.readthedocs.io/en/latest/solvers-boomeramg.html}.

\end{document}
